\documentclass[10pt, twocolumn]{article}

% ─── Codificación y lenguaje ───────────────────────────────────────────────────
\usepackage[utf8]{inputenc}
\usepackage[T1]{fontenc}
\usepackage[spanish]{babel}

% ─── Matemáticas ──────────────────────────────────────────────────────────────
\usepackage{amsmath, amssymb, amsthm}

% ─── Tipografía y geometría ────────────────────────────────────────────────────
\usepackage{lmodern}
\usepackage[margin=1.8cm, top=2cm, bottom=2cm]{geometry}

% ─── Tablas y figuras ─────────────────────────────────────────────────────────
\usepackage{booktabs}
\usepackage{tabularx}
\usepackage{array}
\usepackage{graphicx}
\usepackage{float}

% ─── Color y cajas ────────────────────────────────────────────────────────────
\usepackage{xcolor}
\usepackage{tcolorbox}
\tcbuselibrary{breakable, skins}

% ─── Hipervínculos ────────────────────────────────────────────────────────────
\usepackage[hidelinks]{hyperref}

% ─── Caja de calibración ──────────────────────────────────────────────────────
\newtcolorbox{notacal}{
  colback=yellow!10, colframe=orange!60!black,
  fonttitle=\bfseries\small, title=Nota de calibración,
  breakable, left=4pt, right=4pt, top=3pt, bottom=3pt
}

% ─── Marcado de IA ────────────────────────────────────────────────────────────
\newcommand{\ia}[1]{\textcolor{teal}{\textbf{\{ia\}}~#1}}

% ─── Metadatos ────────────────────────────────────────────────────────────────
\title{\textbf{Modelo Dinámico de Estabilidad Térmica y Rendimiento\\
       en Computadores de Propósito General}}

\author{%
  David Alexander Salazar Villa \and
  Julian David Ramírez Rodríguez \and
  Juan Andrés Gaviria
}

\date{\today}

% ══════════════════════════════════════════════════════════════════════════════
\begin{document}
\maketitle

% ─── Responsabilidades ────────────────────────────────────────────────────────
\noindent\textbf{Responsabilidades del equipo:}
\begin{itemize}
  \item \textbf{David Salazar} --- [completar]
  \item \textbf{Julian Ramírez} --- [completar]
  \item \textbf{Juan Gaviria} --- [completar]
\end{itemize}

% ─── Resumen ──────────────────────────────────────────────────────────────────
\begin{abstract}
Se modela la dinámica del \emph{thermal throttling} en procesadores de
computadores de propósito general bajo cargas sostenidas como un sistema de
tres ecuaciones diferenciales no lineales acopladas que relacionan temperatura,
potencia y throughput efectivo. La función de throttling logística (sigmoide)
introduce una no linealidad suave pero irreducible que impide la solución
analítica y divide el comportamiento del sistema en dos regímenes: aceleración
libre y protección térmica. Se implementan tres métodos numéricos programados
desde cero ---Euler, Euler Mejorado y Runge-Kutta de cuarto orden (RK4)--- y
se realiza un análisis comparativo de error variando el tamaño de paso,
validando contra el solucionador de referencia RK45. Se simulan cuatro
escenarios de intervención ---base, \emph{undervolting}, mejora de disipación y
cambio de pasta térmica--- evaluando en cada uno si el sistema converge a un
equilibrio estable con throughput viable. Los resultados se animan para ilustrar
las trayectorias en el espacio de estados.
\end{abstract}

% ══════════════════════════════════════════════════════════════════════════════
\section{Introducción}

\subsection{Presentación del sistema}

El sistema de estudio es la unidad de procesamiento central (CPU) de un
computador de propósito general operando bajo cargas de trabajo sostenidas de
alto rendimiento. En la arquitectura de Von Neumann, el procesador enfrenta una
tensión permanente entre tres fuerzas en competencia: el planificador del sistema
operativo, que inyecta potencia para maximizar el throughput efectivo; la
generación de calor por efecto Joule, proporcional a esa potencia; y los
mecanismos de protección del firmware, que reducen la energía para prevenir el
daño físico irreversible del silicio.

Esta interacción produce una retroalimentación altamente no lineal. Cuando el
SO sube la potencia, la temperatura sube; cuando la temperatura supera un umbral
crítico, el firmware interviene recortando la potencia; al caer la potencia, el
throughput cae por debajo del objetivo y el SO vuelve a demandar más energía.
Dependiendo de los parámetros físicos del hardware, este ciclo puede converger
a un equilibrio estable o mantenerse como una oscilación incontrolable.

\subsection{Relevancia}

El fenómeno es transversal a cualquier computador de propósito general sometido
a carga sostenida: servidores de cómputo científico, estaciones de trabajo de
renderizado, equipos de procesamiento de datos en tiempo real. El modelo permite
predecir analíticamente si un conjunto de parámetros de hardware produce un
sistema estable o inestable, con aplicación directa en el diseño de políticas de
refrigeración y selección de componentes.

\begin{itemize}
  \item El sistema operativo busca maximizar la productividad asignando la mayor
        potencia posible para maximizar el throughput.
  \item El procesador, por efecto Joule, incrementa su temperatura de forma
        proporcional a esa carga de trabajo.
  \item El firmware reduce la energía de forma progresiva y agresiva para
        prevenir daños irreversibles en la arquitectura del chip.
\end{itemize}

\subsection{Fenómeno de estudio}

El \emph{thermal throttling} es el fenómeno central. El firmware no actúa como
un interruptor binario: implementa el Escalado Dinámico de Frecuencia y Voltaje
(DVFS) ---Intel P-States, AMD Precision Boost--- reduciendo la frecuencia de
forma progresiva y acelerada a medida que la temperatura se aproxima al umbral
crítico. Esto genera oscilaciones severas en cargas sostenidas: cuando el equipo
entra en throttling profundo, el throughput se desploma, y al recuperarse la
temperatura el ciclo reinicia, comprometiendo la integridad del procesamiento.

\subsubsection{Por qué no existe solución analítica}

El sistema \textbf{no admite solución analítica exacta} por dos razones
complementarias.

La primera es \textbf{estructural}: la función de throttling logístico introduce
una no linealidad irreducible en la ecuación de control de potencia. No existe
transformación de variables conocida que linealice simultáneamente los tres
estados acoplados.

La segunda es \textbf{analítica}: las tres ecuaciones están doblemente acopladas
---$P$ aparece como término fuente en $\dot{T}$ y en $\dot{R}$; $R$ aparece en
$\dot{P}$; $T$ aparece en $\dot{P}$ a través de la función logística---. Este
acoplamiento circular impide integrar una ecuación de forma independiente y
sustituir su resultado en las demás.

Por tanto, el uso de métodos numéricos no es una simplificación del problema
sino la \textbf{única metodología válida} para estudiar el comportamiento
completo del sistema.

\subsection{Objetivos}

El objetivo principal es modelar matemáticamente y simular mediante métodos
numéricos la dinámica térmica y de rendimiento de un procesador bajo carga
sostenida, identificando las condiciones bajo las cuales el sistema alcanza
un estado de equilibrio estable.

Objetivos específicos: implementar Euler, Euler Mejorado y RK4 programados
desde cero para un sistema vectorial de tres estados; realizar un análisis
comparativo de convergencia y error variando $h$ validando contra RK45; simular
cuatro escenarios de intervención; y producir animaciones de las trayectorias
en el espacio de estados.

\subsection{Estructura del informe}

La \textbf{Sección~\ref{sec:metodologia}} desglosa el modelo matemático y
justifica las constantes físicas. La \textbf{Sección~\ref{sec:numerica}}
describe la metodología numérica, validación y herramientas. La
\textbf{Sección~\ref{sec:resultados}} presenta el análisis de equilibrio,
escenarios y análisis de error. La \textbf{Sección~\ref{sec:conclusiones}}
sintetiza los hallazgos.

% ══════════════════════════════════════════════════════════════════════════════
\section{Metodología}
\label{sec:metodologia}

\subsection{Modelo matemático}

El fenómeno se modela como un sistema de tres EDOs no lineales acopladas de
primer orden. Las variables de estado son:

\begin{itemize}
  \item $T(t)$: temperatura del procesador [\textdegree C]
  \item $P(t)$: potencia eléctrica asignada [W]
  \item $R(t)$: throughput efectivo (operaciones por segundo) [ops/s]
\end{itemize}

Las tres están acopladas: $P$ determina cuánto se calienta el procesador y
cuán rápido procesa; $T$ activa el throttling que recorta $P$; la caída de
$P$ reduce $R$, lo que impulsa al SO a demandar más $P$.

% ─── Ecuación 1 ───────────────────────────────────────────────────────────────
\subsubsection*{Ecuación 1: Dinámica Térmica}

\begin{equation}
  \frac{dT}{dt} = k_1 P(t) - k_2 \bigl(T(t) - T_{\text{amb}}\bigr)
  \label{eq:termica}
\end{equation}

Aplicación directa del Modelo de Capacitancia Térmica Concentrada
(\emph{Lumped Capacitance Model}): el procesador se trata como un nodo único
con masa térmica $C_{\text{th}}$ y resistencia térmica $R_{\text{th}}$ hacia
el ambiente. El balance energético canónico
$C_{\text{th}} \frac{dT}{dt} = \dot{E}_{\text{in}} - \dot{E}_{\text{out}}$
---energía entrante por Efecto Joule, saliente por Ley de Newton--- dividido
entre $C_{\text{th}}$ produce la forma exacta, con
$k_1 \equiv 1/C_{\text{th}}$ y $k_2 \equiv 1/(R_{\text{th}} C_{\text{th}})$.

\textbf{Término} $k_1 P(t)$: tasa de calentamiento. $k_1$
[\textdegree C\,J$^{-1}$] cuantifica cuántos grados produce cada julio;
depende de la litografía del silicio a través de $\kappa_0$.

\textbf{Término} $-k_2(T - T_{\text{amb}})$: disipación térmica por
convección. $k_2$ [s$^{-1}$] cuantifica la eficiencia del disipador: nula
cuando $T = T_{\text{amb}}$, máxima cuando $T \gg T_{\text{amb}}$.

\smallskip
\noindent\textit{Ref.: Incropera \& DeWitt --- \emph{Fundamentos de
transferencia de calor y masa}, cap.\ Conducción transitoria~\cite{incropera}.}

% ─── Ecuación 2 ───────────────────────────────────────────────────────────────
\subsubsection*{Ecuación 2: Dinámica de Rendimiento}

\begin{equation}
  \frac{dR}{dt} = \gamma P(t) - \delta R(t)
  \label{eq:rendimiento}
\end{equation}

Su forma proviene del modelo de fluidos en teoría de colas
(\emph{Fluid Queueing Model}): $\dot{x} = \lambda(t) - \mu x(t)$, donde
$\lambda(t)$ es la tasa de inyección de trabajo y $\mu x(t)$ la fricción
proporcional a la carga. $R(t)$ actúa como la variable de flujo $x(t)$.

\textbf{Término} $\gamma P(t)$: aceleración computacional. A mayor potencia,
mayor frecuencia de operación y mayor throughput. $\gamma$ mide la eficiencia
con que la energía se convierte en operaciones útiles, determinada por el
paralelismo del hardware (hilos lógicos, frecuencia, IPC efectivo).

\textbf{Término} $-\delta R(t)$: fricción computacional. Al aumentar $R$, se
saturan los buses de comunicación, controladores de memoria y colas de
instrucciones. Según el Modelo Roofline, la carga opera en el régimen
\emph{memory bound}: el cuello de botella es el ancho de banda por el que los
datos llegan al procesador. $\delta$ [s$^{-1}$] es inversamente proporcional
al ancho de banda del enlace más lento en la ruta crítica de datos.

\smallskip
\noindent\textit{Ref.: Kleinrock --- \emph{Queueing Systems, Vol.\,II},
cap.\ Fluid Approximations~\cite{kleinrock}; Williams et al.\ ---
Roofline model~\cite{williams}.}

% ─── Ecuación 3 ───────────────────────────────────────────────────────────────
\subsubsection*{Ecuación 3: Control de Potencia con Throttling Logístico}

\begin{equation}
  \frac{dP}{dt} = \alpha\bigl(R_{\text{obj}} - R(t)\bigr)
                 - \frac{\beta}{1 + e^{-k(T(t) - T_{\text{crit}})}}
  \label{eq:potencia}
\end{equation}

Esta ecuación es la \textbf{lógica de control de potencia} del SO y el
firmware, y contiene el término no lineal que hace al sistema irresoluble
de forma analítica.

\textbf{Término} $\alpha(R_{\text{obj}} - R(t))$: demanda del SO. En la
teoría de control automático corresponde a la acción de un controlador
proporcional (componente P de un PID): la ganancia $\alpha$ actúa sobre el
error $e(t) = R_{\text{obj}} - R(t)$ para ajustar la potencia. Un $\alpha$
alto corresponde al perfil ``Máximo rendimiento''; un $\alpha$ bajo,
a ``Ahorro de batería''.

\textbf{Término logístico} $-\frac{\beta}{1 + e^{-k(T - T_{\text{crit}})}}$:
throttling del firmware. Modela el DVFS mediante una función sigmoide
que produce una transición continua y acelerada alrededor de $T_{\text{crit}}$.
Su comportamiento en los tres regímenes:

\begin{itemize}
  \item $T \ll T_{\text{crit}}$: término $\to 0$. El firmware no interviene.
  \item $T = T_{\text{crit}}$: término $= \beta/2$. Corte al 50\,\% del máximo.
  \item $T \gg T_{\text{crit}}$: término $\to \beta$. Corte máximo de potencia.
\end{itemize}

El parámetro $k$ [\textdegree C$^{-1}$] controla la agresividad de la
transición: $k$ alto produce una curva casi vertical; $k$ bajo, respuesta
gradual. Se fija durante la calibración de la simulación.

\medskip
La elección de la función logística frente a alternativas como
$\max(0,\Delta T)^p$ se justifica por dos criterios:

\smallskip\noindent\textbf{1. Realismo físico --- DVFS.}
Los procesadores modernos (Intel P-States, AMD Precision Boost) reducen la
frecuencia de forma progresiva a medida que la temperatura se aproxima al
límite, sin esperar a superarlo. Brooks y Martonosi~\cite{brooks} documentan
que estos mecanismos imponen caídas de rendimiento asimétricas cuya severidad
aumenta de forma no lineal. La sigmoide modela esta transición continua con
mayor fidelidad que una función que permanece en cero hasta $T_{\text{crit}}$
y luego actúa de golpe.

\smallskip\noindent\textbf{2. Diferenciabilidad $C^\infty$.}
La función logística es infinitamente diferenciable en todo $\mathbb{R}$: no
presenta discontinuidades en ninguna derivada en ningún punto, incluyendo la
vecindad de $T_{\text{crit}}$. Esto garantiza que RK4, al evaluar sus cuatro
pendientes en $[t_n, t_{n+1}]$, nunca encuentre un cambio brusco de curvatura.
El orden de convergencia del método se mantiene globalmente, no solo por tramos.
Nocedal y Wright~\cite{nocedal} establecen que la diferenciabilidad continua
de orden superior es la condición que preserva el orden local de los métodos
de integración explícitos.

La constante $\beta$ [W] acota la potencia máxima de corte del firmware. Un
fabricante conservador usará $\beta$ alto; uno orientado al rendimiento, $\beta$
bajo.

\smallskip
\noindent\textit{Refs.: Ogata~\cite{ogata}; Ames et al.~\cite{ames};
Nocedal \& Wright~\cite{nocedal}; Brooks \& Martonosi~\cite{brooks}.}

% ─── Constantes ───────────────────────────────────────────────────────────────
\subsection{Definición de constantes}

\subsubsection{Variables base de hardware}

Todas son fijas al adquirir el hardware. Los únicos parámetros intervenibles
entre escenarios son $\Phi_{\text{aire}}$ (base refrigerante o limpieza)
y $H_d$ (cambio de pasta térmica).

\begin{table}[H]
\centering
\caption{Variables de hardware del modelo}
\label{tab:hardware}
\small
\begin{tabular}{@{}lll@{}}
\toprule
Símbolo & Unidad & Descripción \\
\midrule
$N_c$                    & ---             & Núcleos físicos \\
$N_h$                    & ---             & Hilos lógicos (\emph{threads}) \\
$f_p$                    & GHz             & Frecuencia de reloj \\
$v_{\text{ram}}$         & MT/s            & Velocidad de la RAM \\
$v_{\text{bus}}$         & GT/s            & Velocidad del bus PCIe \\
$\Phi_{\text{aire}}$     & CFM             & Flujo de aire de ventiladores \\
$H_d$                    & W/\textdegree C & Conductancia del disipador \\
\bottomrule
\end{tabular}
\end{table}

\subsubsection{Factores de escala empíricos}

Cada constante global se expresa como el producto de un \textbf{factor de
escala empírico} (subíndice $0$) y una combinación de variables de hardware.
Los factores agrupan fenómenos microscópicos no modelados individualmente
(resistencia del silicio, geometría del chasis, \emph{overhead} del kernel)
y ajustan las unidades para garantizar la consistencia dimensional.

\medskip\noindent\textbf{Dinámica Térmica:}

\begin{itemize}
  \item $k_1 = \kappa_0 \cdot (f_p \cdot N_c)$
        [\textdegree C\,s$^{-1}$\,W$^{-1}$].
        A mayor frecuencia y más núcleos activos, más transistores conmutan
        y mayor es la disipación de energía como calor (Efecto Joule).
        $\kappa_0$ captura la resistencia térmica de la litografía del
        silicio; equivale a $1/C_{\text{th}}$: $k_1 \equiv 1/C_{\text{th}}$.

  \item $k_2 = \eta_0 \cdot (\Phi_{\text{aire}} \cdot H_d)$ [s$^{-1}$].
        La capacidad de enfriamiento depende del flujo de aire y la
        conductancia del disipador. $\eta_0$ corrige la geometría real del
        chasis (obstrucciones, flujo no laminar, presión ambiental);
        equivale a $1/(R_{\text{th}} C_{\text{th}})$.
\end{itemize}

\noindent\textbf{Dinámica de Rendimiento:}

\begin{itemize}
  \item $\gamma = \gamma_0 \cdot (N_h \cdot f_p)$
        [ops\,s$^{-2}$\,W$^{-1}$].
        Basado en la \emph{Iron Law}: el throughput bruto es proporcional
        a hilos lógicos y frecuencia de reloj. $\gamma_0$ representa el
        IPC efectivo de la arquitectura para la carga de trabajo estudiada.

  \item $\delta = \delta_0 \,/\, \min(v_{\text{ram}},\, v_{\text{bus}})$
        [s$^{-1}$].
        Modelo Roofline (\emph{memory bound}): el cuello de botella es el
        ancho de banda del enlace más lento, ya sea la RAM (MT/s) o el bus
        PCIe (GT/s). $\delta_0$ cuantifica el \emph{overhead} intrínseco
        del sistema operativo.
\end{itemize}

\noindent\textbf{Control de Potencia:}

\begin{itemize}
  \item $\alpha = \alpha_0 \cdot N_h/N_c$
        [W\,s$^{-1}$\,(ops\,s$^{-1}$)$^{-1}$].
        La ratio $N_h/N_c$ (índice de \emph{hyperthreading}) escala la
        agresividad del gobernador de energía. $\alpha_0$ representa la
        sensibilidad del perfil energético seleccionado.

  \item $\beta$ [W].
        Amplitud máxima del corte de potencia del firmware. No depende de
        variables de hardware individuales; refleja la política del
        fabricante: alto para diseños conservadores, bajo para orientados
        al rendimiento.

  \item $k$ [\textdegree C$^{-1}$].
        Agresividad de la transición logística. Se fija durante la
        calibración empírica del comportamiento DVFS del procesador.
\end{itemize}

% ─── Análisis dimensional ─────────────────────────────────────────────────────
\subsection{Análisis dimensional}

Las unidades de las variables de estado son:
$T(t)$ [\textdegree C], $R(t)$ [ops/s], $P(t)$ [W].

\medskip\noindent\textbf{Ecuación térmica}
($dT/dt$ en \textdegree C\,s$^{-1}$):
\[
\underbrace{k_1}_{\left[\frac{^\circ\text{C}}{\text{W\,s}}\right]}
\underbrace{P}_{\left[\text{W}\right]}
= \left[\frac{^\circ\text{C}}{\text{s}}\right], \qquad
\underbrace{k_2}_{\left[\text{s}^{-1}\right]}
\underbrace{(T - T_{\text{amb}})}_{\left[^\circ\text{C}\right]}
= \left[\frac{^\circ\text{C}}{\text{s}}\right]
\]

\noindent\textbf{Ecuación de rendimiento}
($dR/dt$ en ops\,s$^{-2}$):
\[
\underbrace{\gamma}_{\left[\frac{\text{ops}}{\text{W\,s}^2}\right]}
\underbrace{P}_{\left[\text{W}\right]}
= \left[\frac{\text{ops}}{\text{s}^2}\right], \qquad
\underbrace{\delta}_{\left[\text{s}^{-1}\right]}
\underbrace{R}_{\left[\frac{\text{ops}}{\text{s}}\right]}
= \left[\frac{\text{ops}}{\text{s}^2}\right]
\]

\noindent\textbf{Ecuación de potencia} ($dP/dt$ en W\,s$^{-1}$):
\[
\underbrace{\alpha}_{\left[\frac{\text{W\,s}^{-1}}{\text{ops\,s}^{-1}}\right]}
\underbrace{(R_{\text{obj}}-R)}_{\left[\frac{\text{ops}}{\text{s}}\right]}
= \left[\frac{\text{W}}{\text{s}}\right]
\]
\[
\frac{
  \underbrace{\beta}_{\left[\text{W}\right]}
}{
  1 + e^{
    -\underbrace{k}_{\left[^\circ\text{C}^{-1}\right]}
    \underbrace{(T-T_{\text{crit}})}_{\left[^\circ\text{C}\right]}
  }
}
= \left[\text{W}\right]
\;\Rightarrow\;
\frac{d}{dt}(\cdots) = \left[\frac{\text{W}}{\text{s}}\right] \checkmark
\]

% ══════════════════════════════════════════════════════════════════════════════
\section{Metodología Numérica}
\label{sec:numerica}

\subsection{Implementación de los tres métodos}

El sistema vectorial $\dot{\mathbf{x}} = \mathbf{f}(\mathbf{x})$ con
$\mathbf{x} = (T, R, P)^\top$ se resuelve con tres métodos programados
desde cero.

\medskip\noindent\textbf{Euler explícito:}
\begin{equation}
  \mathbf{x}_{n+1} = \mathbf{x}_n + h\,\mathbf{f}(\mathbf{x}_n)
\end{equation}

\noindent\textbf{Euler Mejorado (Heun):}
\begin{align}
  \mathbf{x}^* &= \mathbf{x}_n + h\,\mathbf{f}(\mathbf{x}_n) \notag \\
  \mathbf{x}_{n+1} &= \mathbf{x}_n
    + \frac{h}{2}\bigl(\mathbf{f}(\mathbf{x}_n) + \mathbf{f}(\mathbf{x}^*)\bigr)
\end{align}

\noindent\textbf{Runge-Kutta de cuarto orden (RK4):}
\begin{align*}
  \mathbf{k}_1 &= \mathbf{f}(\mathbf{x}_n) \\
  \mathbf{k}_2 &= \mathbf{f}\!\left(\mathbf{x}_n + \tfrac{h}{2}\mathbf{k}_1\right) \\
  \mathbf{k}_3 &= \mathbf{f}\!\left(\mathbf{x}_n + \tfrac{h}{2}\mathbf{k}_2\right) \\
  \mathbf{k}_4 &= \mathbf{f}(\mathbf{x}_n + h\,\mathbf{k}_3)
\end{align*}
\begin{equation}
  \mathbf{x}_{n+1} = \mathbf{x}_n
    + \frac{h}{6}(\mathbf{k}_1 + 2\mathbf{k}_2 + 2\mathbf{k}_3 + \mathbf{k}_4)
  \label{eq:rk4}
\end{equation}

\noindent\textbf{Condición de frontera programática sobre $P(t)$.}
La ODE permite teóricamente que $P(t)$ sea negativa bajo throttling máximo
sostenido, lo que carece de sentido físico. Al finalizar cada paso de
integración se aplica la corrección:
\begin{equation}
  P_{n+1} \leftarrow \max\!\bigl(P_{\min},\; P_{n+1}\bigr),
  \quad P_{\min} = 5\;\text{W}
\end{equation}
Esta corrección \textbf{no modifica la ODE}; es una regla lógica en el
código que descarta resultados físicamente imposibles, garantizando que
el modelo permanezca en un dominio admisible.

\subsection{Inestabilidad inducida por gradientes extremos}

RK4 es un integrador \emph{explícito}: en cada paso calcula cuatro
pendientes locales y las combina para extrapolar el estado. Funciona bien
cuando las pendientes cambian suavemente en $[t_n, t_{n+1}]$.

La función logística tiene derivada máxima en $T = T_{\text{crit}}$, donde
vale $k\beta/4$. Para valores altos de $k$ (transición agresiva), esta
derivada puede ser grande y el cambio de pendiente dentro del intervalo
puede ser significativo si $h$ es demasiado grande. En ese caso, RK4
sobreestima o subestima el corte de potencia, produciendo que $P_{n+1}$
se aleje del valor real y propague error al siguiente paso.

La ventaja de la sigmoide frente a otras funciones de penalización es que,
al ser $C^\infty$, la degradación del orden de convergencia de RK4 es
gradual y proporcional a $k$ y $h$, no abrupta como ocurriría con una
función no diferenciable. La solución es reducir $h$ hasta que la función
logística no cambie sustancialmente dentro de un paso, condición que se
verifica empíricamente en el análisis de convergencia.

\subsection{Análisis de error y convergencia}

Para cada $h \in \{1.0,\; 0.5,\; 0.1,\; 0.05,\; 0.01\}$ segundos, se
calcula el error relativo global respecto a RK45:

\begin{equation}
  E_{\text{rel}}(h) = \max_{t \in [0, T_{\text{sim}}]}
  \frac{\|\mathbf{x}_{\text{aprox}}(t;\,h) - \mathbf{x}_{\text{RK45}}(t)\|}
       {\|\mathbf{x}_{\text{RK45}}(t)\|}
\end{equation}

Se espera convergencia de orden 1 para Euler ($E \sim h$), orden 2 para
Euler Mejorado ($E \sim h^2$) y orden 4 para RK4 ($E \sim h^4$). La
gráfica log-log de $E_{\text{rel}}$ vs $h$ debe mostrar estas pendientes
características.

\subsection{Validación}

La validación se realiza comparando contra
\texttt{scipy.integrate.solve\_ivp} con método RK45 de paso adaptativo.
La implementación se considera correcta cuando $E_{\text{rel}}$ disminuye
con el orden esperado al reducir $h$, y cuando los tres métodos producen
los mismos puntos de equilibrio $(T^*, R^*, P^*)$ y el mismo comportamiento
cualitativo en cada escenario.

\subsection{Herramientas}

Implementación en Python~3.x: \texttt{NumPy} (operaciones vectoriales),
\texttt{SciPy} (\texttt{solve\_ivp}/RK45, referencia), \texttt{Matplotlib}
(visualización estática) y \texttt{Manim} (animaciones). El código de los
tres métodos está programado directamente por el equipo.

\subsection{Proceso de programación, simulación y animación}

Se define primero la función vectorial $\mathbf{f}(\mathbf{x})$ que agrupa
las tres ecuaciones, aplicando la condición de frontera sobre $P_{\min}$ al
final de cada evaluación. Sobre esta función se construyen los tres
integradores con la misma interfaz, lo que permite intercambiarlos
directamente en el bucle de simulación.

Los cuatro escenarios se ejecutan con el mismo horizonte temporal y
condiciones iniciales, variando únicamente los parámetros de hardware
correspondientes. Los resultados se grafican como series temporales de
$T$, $R$ y $P$, y como trayectorias en el espacio de fases $(T, R, P)$.
Las animaciones con Manim muestran la evolución en tiempo real, resaltando
la activación del throttling, la entrada al régimen oscilatorio y la
convergencia al equilibrio.

\subsection{Uso de IA}

\begin{itemize}
  \item \ia{Consulta de conceptos teóricos sobre el Modelo Roofline,
        Capacitancia Concentrada y funciones de barrera de control.}
  \item \ia{Identificación de referencias bibliográficas relevantes.}
  \item \ia{Revisión de errores en la verificación dimensional de las
        constantes tras un análisis propio previo.}
  \item \ia{Sugerencia de condiciones iniciales y rangos de parámetros
        interesantes, explorada después de forma propia.}
\end{itemize}

No se empleó IA para: generar el código, escribir el informe o la
presentación, producir imágenes o animaciones, ni tomar decisiones de
modelado o análisis.

% ══════════════════════════════════════════════════════════════════════════════
\section{Resultados y Discusión}
\label{sec:resultados}

\subsection{Condiciones iniciales y parámetros}

Las condiciones iniciales representan el estado del equipo en $t = 0$,
justo antes de lanzar la carga de trabajo (equipo en reposo):

\begin{table}[H]
\centering
\caption{Condiciones iniciales}
\label{tab:ci}
\small
\begin{tabular}{@{}lcl@{}}
\toprule
Variable & Valor & Significado físico \\
\midrule
$T(0)$ & 45\,\textdegree C & Temperatura de reposo con carga base del SO \\
$P(0)$ & 15\,W             & Potencia base en inactividad \\
$R(0)$ & 0\,ops/s          & Procesador sin carga activa \\
\bottomrule
\end{tabular}
\end{table}

\begin{table}[H]
\centering
\caption{Parámetros fijos del entorno}
\label{tab:params}
\small
\begin{tabular}{@{}lcl@{}}
\toprule
Parámetro & Valor & Descripción \\
\midrule
$T_{\text{amb}}$  & 25\,\textdegree C     & Temperatura ambiente \\
$T_{\text{crit}}$ & 90\,\textdegree C     & Umbral del throttling logístico \\
$R_{\text{obj}}$  & 20\,ops/s (norm.)     & Throughput objetivo del gobernador \\
$R_{\text{min}}$  & 5\,ops/s (norm.)      & Umbral mínimo de viabilidad \\
$P_{\text{min}}$  & 5\,W                  & Potencia mínima física \\
\bottomrule
\end{tabular}
\end{table}

\begin{table}[H]
\centering
\caption{Hardware de referencia (procesador 8 núcleos, 3.5\,GHz)}
\label{tab:hw}
\small
\begin{tabular}{@{}llcc@{}}
\toprule
Variable & Símbolo & Valor & Unidad \\
\midrule
Núcleos físicos      & $N_c$                   & 8    & ---  \\
Hilos lógicos        & $N_h$                   & 16   & ---  \\
Frecuencia           & $f_p$                   & 3.5  & GHz  \\
Velocidad RAM        & $v_{\text{ram}}$        & 4800 & MT/s \\
Velocidad bus PCIe   & $v_{\text{bus}}$        & 16   & GT/s \\
Flujo de aire        & $\Phi_{\text{aire}}$    & 2.5  & CFM  \\
Conductancia disip.  & $H_d$                   & 5.0  & W/\textdegree C \\
\bottomrule
\end{tabular}
\end{table}

\begin{table}[H]
\centering
\caption{Factores de escala empíricos}
\label{tab:escala}
\small
\begin{tabular}{@{}llcl@{}}
\toprule
Factor & Símbolo & Valor & Justificación \\
\midrule
Ineficiencia térmica & $\kappa_0$ & $0.002$           & Litografía 7\,nm \\
Convección           & $\eta_0$   & $0.004$           & Chasis parcialmente obstruido \\
Eficiencia IPC       & $\gamma_0$ & $0.015$           & Cómputo paralelo \\
Latencia intrínseca  & $\delta_0$ & $1.2\times10^{8}$ & \emph{Overhead} estándar del SO \\
Ganancia gobernador  & $\alpha_0$ & $0.8$             & Perfil Máximo Rendimiento \\
Corte máx.\ firmware & $\beta$    & $12$\,W           & Protección moderada \\
Agresividad logística& $k$        & $0.5$             & DVFS típico [\textdegree C$^{-1}$] \\
\bottomrule
\end{tabular}
\end{table}

Con estos valores, los parámetros del modelo resultan:

\begin{align}
  k_1    &= 0.002 \times (3.5 \times 8)
           = 0.056\;\text{\textdegree C\,s}^{-1}\text{W}^{-1} \notag\\
  k_2    &= 0.004 \times (2.5 \times 5.0)
           = 0.050\;\text{s}^{-1} \notag\\
  \gamma &= 0.015 \times (16 \times 3.5)
           = 0.840\;\text{ops\,s}^{-2}\text{W}^{-1} \notag\\
  \delta &= \frac{1.2\times10^{8}}{\min(4800,\,16)}
           = \frac{1.2\times10^{8}}{16}
           = 7.5\times10^{6}\;\text{s}^{-1} \notag\\
  \alpha &= 0.8 \times \frac{16}{8}
           = 1.6\;\text{W\,s}^{-1}(\text{ops\,s}^{-1})^{-1} \notag
\end{align}

\begin{notacal}
Los factores $\gamma_0$ y $\delta_0$ deben ajustarse iterativamente hasta
que $P^* = \delta \cdot R_{\text{obj}}/\gamma$ caiga en el rango físico
esperado (15--45\,W bajo carga). Este proceso de calibración es parte de la
metodología y debe documentarse en el informe.
\end{notacal}

\subsection{Análisis de puntos de equilibrio}

\subsubsection{Equilibrio sin throttling activo}

Igualando las tres derivadas a cero y asumiendo que $T^*$ está
suficientemente alejada de $T_{\text{crit}}$ (throttling despreciable):

\begin{align*}
  \frac{dT}{dt}=0 &\implies T^* = T_{\text{amb}} + \frac{k_1}{k_2} P^* \\
  \frac{dR}{dt}=0 &\implies R^* = \frac{\gamma}{\delta} P^* \\
  \frac{dP}{dt}=0 &\implies R^* = R_{\text{obj}}
\end{align*}

De la tercera se obtiene $R^* = R_{\text{obj}}$. Sustituyendo:

\begin{equation}
  P^* = \frac{\delta \cdot R_{\text{obj}}}{\gamma}, \qquad
  T^* = T_{\text{amb}} + \frac{k_1 \delta \cdot R_{\text{obj}}}{k_2 \gamma}
\end{equation}

La \textbf{condición de viabilidad térmica} es:

\begin{equation}
  T^* < T_{\text{crit}}
  \iff
  \frac{k_1 \delta}{k_2 \gamma} < \frac{T_{\text{crit}} - T_{\text{amb}}}{R_{\text{obj}}}
  \label{eq:viabilidad}
\end{equation}

Cuando se cumple, el sistema puede alcanzar un equilibrio estable con
throughput pleno. Cuando no, entra en régimen oscilatorio.

\subsubsection{Estabilidad por retroalimentación negativa}

\noindent\textbf{Perturbación térmica.}
Si $T$ sube a $T^* + \varepsilon$ con $\varepsilon > 0$:

\begin{equation*}
  \frac{dT}{dt}\bigg|_{T^*+\varepsilon}
  = \underbrace{k_1 P^* - k_2(T^*-T_{\text{amb}})}_{=\,0}
  - k_2\varepsilon = -k_2\varepsilon < 0
\end{equation*}

La derivada es negativa: el sistema enfría activamente hacia $T^*$.
$k_2 > 0$ actúa como fuerza de amortiguación proporcional al desvío
---retroalimentación negativa pura.

\noindent\textbf{Perturbación de rendimiento.}
Si $R$ cae bajo $R^*$, el término $\alpha(R_{\text{obj}} - R)$ se vuelve
positivo, aumentando $\dot{P}$, lo que empuja $\dot{R} = \gamma P - \delta R$
hacia valores positivos y recupera el throughput. El término $-\delta R(t)$
actúa como fricción que impide sobrepasar $R^*$.

La estabilidad requiere que se satisfaga la
condición~(\ref{eq:viabilidad}).

\subsubsection{Régimen oscilatorio}

Cuando la condición de viabilidad no se cumple:

\begin{enumerate}
  \item El SO aumenta $P$ buscando alcanzar $R_{\text{obj}}$.
  \item $T$ sube y activa progresivamente el throttling logístico.
  \item El firmware reduce $P$ de forma acelerada.
  \item $T$ cae al perder la fuente de calor.
  \item El SO intenta recuperar $P$ y el ciclo reinicia.
\end{enumerate}

El equipo es inviable para cargas sostenidas: $R$ oscila sin converger
a un estado estable con $R \geq R_{\min}$.

\subsection{Cuatro escenarios de intervención}

\subsubsection*{Escenario 0 --- Base (control)}
Hardware sin modificaciones. Parámetros de la Tabla~\ref{tab:hw}. Puede o
no satisfacer la condición~(\ref{eq:viabilidad}). Punto de referencia para
variaciones porcentuales.

\subsubsection*{Escenario 1 --- Undervolting (perfil conservador)}
\textbf{Modificación:} Reducción de $\alpha_0$.\\
\textbf{Efecto esperado:} El SO sube $P$ más lentamente, alejando $T$ de
$T_{\text{crit}}$. Puede estabilizar el sistema a costa de un $R^*$
inferior.\\
\textbf{Comparación:} ¿Alcanza $R^* \geq R_{\min}$?

\subsubsection*{Escenario 2 --- Hardware ligero (base refrigerante)}
\textbf{Modificación:} Aumento de $\Phi_{\text{aire}}$ $\to$ $k_2$ sube.\\
\textbf{Efecto esperado:} Mayor disipación desplaza $T^*$ hacia abajo sin
sacrificar throughput.\\
\textbf{Comparación:} Variación porcentual de $T^*$ y frecuencia de
oscilación vs.\ escenario base.

\subsubsection*{Escenario 3 --- Hardware profundo (metal líquido)}
\textbf{Modificación:} Aumento de $H_d$ $\to$ $k_2$ sube sustancialmente.\\
\textbf{Efecto esperado:} Si el nuevo $k_2$ cumple~(\ref{eq:viabilidad}),
el sistema converge al equilibrio con throughput pleno.\\
\textbf{Comparación:} Tiempo de convergencia y temperatura máxima durante
la transición.

\subsection{Comparación entre métodos y tamaños de paso}

% PENDIENTE: completar con resultados numéricos.
% Incluir tabla de E_rel(h) para Euler, Euler Mejorado y RK4.
% Incluir gráfica log-log de E vs h con pendientes de orden 1, 2 y 4.
\textit{[Sección por completar con resultados de la simulación.]}

% ══════════════════════════════════════════════════════════════════════════════
\section{Conclusiones}
\label{sec:conclusiones}

% PENDIENTE: completar tras obtener resultados de simulación.
\textit{[Sección por completar.]}

% ══════════════════════════════════════════════════════════════════════════════
\section*{Repositorio}

Código disponible en: \url{https://github.com/...}

% ══════════════════════════════════════════════════════════════════════════════
\begin{thebibliography}{99}

\bibitem{incropera}
Incropera, F.\,P.\ \& DeWitt, D.\,P.\ (2002).
\emph{Fundamentos de transferencia de calor y masa}.
Wiley. Cap.\ Conducción transitoria / Sistemas de capacitancia concentrada.

\bibitem{kleinrock}
Kleinrock, L.\ (1976).
\emph{Queueing Systems, Volume II: Computer Applications}.
Wiley-Interscience. Cap.\ Fluid Approximations.

\bibitem{ogata}
Ogata, K.\ (2010).
\emph{Ingeniería de control moderna}.
Pearson. Cap.\ Acciones básicas de control.

\bibitem{ames}
Ames, A.\,D., Coogan, S., Egerstedt, M., Notomista, G., Sreenath, K.,
\& Tabuada, P.\ (2019).
Control Barrier Functions: Theory and Applications.
\emph{18th European Control Conference (ECC)}.

\bibitem{nocedal}
Nocedal, J.\ \& Wright, S.\,J.\ (2006).
\emph{Numerical Optimization} (2nd ed.).
Springer.

\bibitem{rao}
Rao, S.\,S.\ (2019).
\emph{Engineering Optimization: Theory and Practice} (5th ed.).
Wiley.

\bibitem{brooks}
Brooks, D.\ \& Martonosi, M.\ (2001).
Dynamic Thermal Management for High-Performance Microprocessors.
\emph{Proceedings of the 7th International Symposium on High-Performance
Computer Architecture (HPCA)}.

\bibitem{hennessy}
Hennessy, J.\,L.\ \& Patterson, D.\,A.\ (2019).
\emph{Computer Architecture: A Quantitative Approach} (6th ed.).
Morgan Kaufmann.

\bibitem{williams}
Williams, S., Waterman, A., \& Patterson, D.\ (2009).
Roofline: An Insightful Visual Performance Model for Multicore Architectures.
\emph{Communications of the ACM}, 52(4), 65--76.

\end{thebibliography}

\end{document}
